\documentclass[11pt,a4paper]{article}

%---------------------------------------------------------
%  Packages
%---------------------------------------------------------
\usepackage[utf8]{inputenc}
\usepackage[T1]{fontenc}
\usepackage{geometry}      % Margins
\geometry{margin=1in}
\usepackage{lmodern}
\usepackage{microtype}
\usepackage{hyperref}
\usepackage{enumitem}
\usepackage{amsmath,amssymb}

% Hyperref setup
\hypersetup{
    colorlinks=true,
    linkcolor=blue,
    citecolor=blue,
    urlcolor=blue
}

%---------------------------------------------------------
%  Document
%---------------------------------------------------------
\begin{document}

\section*{Relevant Foundational Work for Predictive Oncology}

\subsection*{Paper~1: \textit{Zhong J} \textit{et~al.}, \textit{EBioMedicine}, 2019\\(Multimodal Fusion for Cancer Prognosis)}

While its primary goal was predicting cancer recurrence, this study provides a powerful methodological template for data integration.\cite{Zhong2019} The authors developed a framework that fuses three distinct data types to predict patient outcomes:

\begin{itemize}[itemsep=2pt]
  \item \textbf{Genomic Data:} Gene expression profiles.
  \item \textbf{Radiomic Data:} Quantitative features extracted from medical images.
  \item \textbf{Clinical Data:} Standard patient demographic and clinical information.
\end{itemize}

The core algorithm is a \textbf{Support Vector Machine (SVM)}, which demonstrated remarkable predictive accuracy. The fused, multimodal model achieved an Area Under the Curve (AUC) of up to \textbf{0.95}, significantly outperforming models trained on any single data modality.

\textbf{Significance.}
This work is crucial because it empirically validates the core premise of our project: that integrating disparate data sources (genomics, imaging, clinical) yields a more powerful predictive signal than any single source alone. It provides a strong precedent for our multimodal fusion strategy. The code for this paper is not explicitly provided.

%---------------------------------------------------------
\subsection*{Paper~2: \textit{Ardila D} \textit{et~al.}, \textit{Nature Medicine}, 2019\\(Deep Learning for Lung Cancer Diagnosis from CT)}
%---------------------------------------------------------

This landmark study from Google Health demonstrates the state-of-the-art in applying deep learning to medical imaging for primary cancer diagnosis.\cite{Ardila2019} The research focused on predicting the malignancy of lung nodules from low-dose chest computed tomography (LDCT) scans.

\begin{enumerate}[itemsep=2pt]
    \item \textbf{Data Used:} The model was trained on a massive dataset of \textbf{42,290} LDCT scans from clinical trial cohorts.
    \item \textbf{AI Algorithm:} The architecture is a \textbf{3D Convolutional Neural Network (CNN)}, which allows the model to learn from the full volumetric context of the CT scan, including nodule morphology and its surroundings.
\end{enumerate}

The model's performance was compared against board-certified radiologists. The deep learning system achieved an \textbf{AUC of 0.94}, demonstrating performance on par with or exceeding that of human experts.

\textbf{Significance.}
This paper is essential reading as it establishes a high-water mark for deep learning performance in thoracic CT imaging. It provides a clear blueprint for an end-to-end deep learning approach to cancer screening and serves as a key performance benchmark that modern architectures in this domain should aspire to meet or surpass. While the dataset is not public, the methodology is well-described.

%---------------------------------------------------------
%  References  (Placeholders for BibTeX)
%---------------------------------------------------------
\begin{thebibliography}{9}
\bibitem{Zhong2019}
Zhong~J, \textit{et~al}.
A deep learning-based multi-modal fusion framework for prediction of cancer prognosis.
\textit{EBioMedicine}. 2019;49:151-160.

\bibitem{Ardila2019}
Ardila~D, \textit{et~al}.
End-to-end lung cancer screening with three-dimensional deep learning on low-dose chest computed tomography.
\textit{Nature Medicine}. 2019;25:954–961.
\end{thebibliography}

\end{document}